% FISH paper — outline scaffold
% Drop into Overleaf as `paper.tex` (single-file, no extra deps).
% Compiles with pdflatex (default Overleaf engine).
% Double-column conference layout via IEEEtran.

\documentclass[conference]{IEEEtran}

% Hyphenation patterns; Overleaf has these by default
\IEEEoverridecommandlockouts

\usepackage{cite}
\usepackage{amsmath,amssymb,amsfonts}
\usepackage{algorithm}
\usepackage{algpseudocode}
\usepackage{graphicx}
\usepackage{textcomp}
\usepackage{xcolor}
\usepackage[colorlinks=true,linkcolor=black,citecolor=blue,urlcolor=blue]{hyperref}
\usepackage{booktabs}
\usepackage[normalem]{ulem}

% Quick TODO marker — visible in PDF
\newcommand{\todoit}[1]{\textcolor{red}{\textbf{[TODO: #1]}}}

\begin{document}

\title{FISH: A Profiling Framework for ROS\,2 Workloads with GPU Acceleration}

\author{\IEEEauthorblockN{Fatih Ta\c{s}yaran}
\IEEEauthorblockA{Sabanc{\i} University \\
\texttt{fatihtasyaran@sabanciuniv.edu}}}

\maketitle

\begin{abstract}
\todoit{abstract — 150--200 words. State the gap: ROS\,2 task models
ignore the GPU side; existing tracers stop at the host API. Summarise
FISH: a four-layer extraction stack (period $\to$ shape $\to$ phase
$\to$ stream-bridge attribution) that produces per-callback typed DAGs
and a process-level task graph from a single trace. Headline numbers:
13$\times$ trace volume reduction via whitelist; \textgreater 99\,\%
attribution coverage on Autoware pointcloud; conditional GPU-DAG
extraction with shared backbone + variant branches.}
\end{abstract}

% --- Introduction ---------------------------------------------------------
\section{Introduction}
\todoit{motivate: ROS\,2 + CUDA workloads in autonomous driving;
existing trace tooling treats the GPU as opaque.}

\subsection{Contributions}
\begin{itemize}
  \item A four-layer trace-extraction stack for ROS\,2+CUDA workloads,
        composing period detection, shape detection, phase classification,
        and stream-bridge attribution.
  \item A typed DAG model (CPU/SM/CE nodes with launch / stream\_fifo /
        cpu\_order edges) derived from CUPTI activity and ROS\,2 tracing.
  \item Skeleton extraction via longest-common-subsequence over per-shape
        DAGs, surfacing a common backbone plus variant branches as
        a conditional graph.
  \item A rule-based per-shape phase classifier
        (\textit{dominant}, \textit{init}, \textit{shutdown}, \textit{anomaly},
        \textit{variant}) with robust outlier detection.
  \item An empirical study on Autoware and a custom drone stack (AAS)
        showing host-mediated stream serialisation dominates real-world
        workloads.
\end{itemize}

% --- Background -----------------------------------------------------------
\section{Background}

\subsection{ROS\,2}
\todoit{node--executor--callback model, callback groups
(MutuallyExclusive vs.\ Reentrant), composable nodes /
component containers.}

\subsection{CUDA}
\todoit{stream model, kernel/memcpy/memset launches, runtime vs.\ driver
API, CUPTI activity records, cudaEvent vs.\ cudaStreamSynchronize
semantics.}

% --- Related Work ---------------------------------------------------------
\section{Related Work}

\subsection{Real-time task models}
\todoit{event-driven DAG, sporadic / periodic, EDD heterogeneous
extensions, response-time analysis on graphs.}

\subsection{Task-extraction tools}
\todoit{ros2\_tracing, CLARA, Pipeline-aware analyses, prior work on
inferring task structure from traces.}

\subsection{Introspection tools}
\todoit{Nsight Systems / Compute, Tracy, Perfetto, ROS-side tracers.
Highlight the gap each one leaves.}

\subsection{Cross-comparison table}
\todoit{Table: tool $\times$ \{ROS-aware, GPU-aware, callback boundaries,
typed DAG, period detection, anomaly detection, paper-ready figures\}.
FISH should be the only entry with all checkmarks.}

% --- Methodology ----------------------------------------------------------
\section{Methodology}

\subsection{Integrating FISH into the environment}
\todoit{transparent \texttt{ros2} wrapper, LTTng tracepoint overlay,
\texttt{nsys} relaunch via launch\_wrap, settings.ini knobs.}

\subsection{Used and added tracepoints}
\todoit{stock 14 events used, 17 FISH-custom events added (scheduler
introspection, callback-link, service client, rclpy parity, action
server). Volume + selection rationale; cite the
\texttt{fish\_events.txt} whitelist (45$\to$31).}

\subsection{FISH stable signal: reproducibility}
\todoit{daemon dual-poll stability (component\_list +
executor\_started); seen-non-empty / seen-non-zero guards; auto-stop
on rosbag end. Why it matters for fair comparison runs.}

% --- Observation section --------------------------------------------------
\section{Expected vs.\ Observed Behaviour of ROS\,2 and Real Applications}
\todoit{motivating observations: framework init dominates t=0;
callback boundary $\neq$ task boundary; GPU side has its own
sync regime; multi-stream serialisation is host-mediated in
practice.}

% --- FISH Model -----------------------------------------------------------
\section{The FISH Model}

\subsection{Mathematical model}
\todoit{layers $L_{-1}$ container, $L_0$ executor, $L_1$ node,
$L_2$ entity, $L_3$ callback. Vertex set $V$, edge set $E$, expansion
operator $\mathcal{E}: V \to 2^V$.}

\subsection{Exact chains from tracepoints}
\todoit{per concept (process--executor, entity, callback, action,
publish-link), the chain of LTTng events that produce each vertex / edge.
Tabulate.}

\subsection{Modelling concepts}
\textbf{Process--executor.} \todoit{process tree, threads, executors,
callback groups.}

\textbf{Entities.} \todoit{publishers, subscriptions, services,
clients, timers, action servers/clients.}

\textbf{Actions.} \todoit{action goal/cancel/result execution.}

\textbf{CPU--GPU interactions.} \todoit{callback windows on CPU side,
typed DAG on GPU side, stream-bridge attribution joining them.}

\textbf{Communication / dependency patterns.} \todoit{publish-link,
client--server, intra-callback chains.}

\subsection{Modelling CPU--GPU interaction}

\subsubsection{CPU--GPU dependencies}
\todoit{three-flag model: stream-id disambiguator, INIT non-prefix,
oort spectrum, callback $\neq$ task. Cite
\texttt{notes/three\_flags\_gpu\_model.txt}.}

\subsubsection{Resolving multi-stream dependencies}
\todoit{deterministic per-stream wait counts; heuristic src$\to$dst
edges from time-proximity; finding: real workloads use
\texttt{cudaStreamSynchronize}, not \texttt{cudaStreamWaitEvent}, so
deterministic edge inference is unattainable for them. Reference
\texttt{notes/host\_mediated\_stream\_sync.txt}.}

\subsubsection{Conditionality on the GPU side: phase and shape detection}
\todoit{shape signature = SHA1((pe,name) sequence $\oplus$ sorted edges).
Phase classes: dominant / init / shutdown / anomaly / variant. Robust
anomaly via mean+$\sigma$ OR median+MAD modified $z$ OR trivial-small
total. Init scoring uses CUDA-init signal tokens (\texttt{cuInit},
\texttt{cuLibraryLoadData}, \ldots).}

\subsubsection{Minimising GPU task graphs}
\todoit{LCS subseq + LCS substring across invocations $\to$
skeleton; per-position branch trie with absolute (entry) and
conditional (continue) probability labels.}

% --- Algorithms -----------------------------------------------------------
\section{Algorithms}

\begin{algorithm}[ht]
\caption{FISH graph generation}
\label{alg:fish-graph}
\begin{algorithmic}[1]
\Require ROS\,2 LTTng trace, process snapshots
\Ensure FISH graph $G = (V, E)$
\todoit{vertex discovery (executor/node/entity/callback) +
horizontal edges (entity / external / propagated /
aggregated).}
\end{algorithmic}
\end{algorithm}

\begin{algorithm}[ht]
\caption{GPU-DAG generation}
\label{alg:gpu-dag}
\begin{algorithmic}[1]
\Require Per-pid CUPTI activity tables, callback windows
\Ensure Per-callback typed DAG list (CPU/SM/CE)
\todoit{stream-bridge attribution then \texttt{\_build\_invocation\_dag}
per window; signature hash; group invocations by shape.}
\end{algorithmic}
\end{algorithm}

\begin{algorithm}[ht]
\caption{GPU task minimisation (skeleton extraction)}
\label{alg:skeleton}
\begin{algorithmic}[1]
\Require A callback's gpu\_dags
\Ensure Skeleton sequence + per-position branch trie
\todoit{multi-sequence LCS over CPU (or all-PE) tokens; align each
invocation to skeleton; aggregate insertions into a trie with class
counts.}
\end{algorithmic}
\end{algorithm}

% --- Taskset extraction ---------------------------------------------------
\section{Taskset Extraction from FISH}

\subsection{EDD-like heterogeneous delay and dependency}
\todoit{per task $\tau_i$: period $T_i$, CPU cost $C^{cpu}_i$,
SM cost $C^{sm}_i$, CE cost $C^{ce}_i$, deadline $D_i$, jitter
$J_i$, dependency edges $\rightarrow$.}

\begin{algorithm}[ht]
\caption{FISH to (modified) EDD}
\label{alg:fish-edd}
\begin{algorithmic}[1]
\Require FISH graph $G$, per-callback skeletons + periods
\Ensure Heterogeneous taskset $\Gamma$
\todoit{walk $G$, fold callback chains into tasks, attach skeleton-
derived costs / period statistics.}
\end{algorithmic}
\end{algorithm}

% --- Verification ---------------------------------------------------------
\section{Verification}

\subsection{Verification of the FISH implementation}
Two routes for ground truth:
\begin{itemize}
  \item Manually inspected small applications (1--3 callbacks).
  \item Procedurally generated workloads with 10s--1000s of nodes /
        callbacks; ground-truth taskset known by construction.
\end{itemize}

% --- Results --------------------------------------------------------------
\section{Results}

\subsection{Accuracy in component discovery}
\todoit{expect 100\,\% on procedurally generated and on AAS / Autoware
spot checks.}

\subsection{Overhead}
\todoit{LTTng + nsys overhead vs.\ baseline; per-callback cost in
$\mu$s; impact on autoware response-time.}

\subsection{Trace data volume}
\todoit{whitelist effectiveness: 36.5M $\to$ 2.3M events
(\textasciitilde 16$\times$); 1.3\,GB $\to$ 93\,MB; 17\,min $\to$
1.5\,min ingest. Reference
\texttt{notes/presentation\_tables\_en.html} Table~3.}

\subsection{Autoware: graph and taskset properties}
\todoit{number of executors, nodes, entities; client-server response
turnaround; pointcloud frame latency.}

\subsection{AAS and other workloads}
\todoit{single-drone mission (takeoff/orbit/land); contrast with
Autoware in callback density and GPU pattern.}

\subsection{GPU task phases and shapes}
\todoit{shape embeddings (UMAP), bimodal cluster of small / large
shapes; F\#1696 case study --- 25 sigs from 45 invocations, 23-token
skeleton common to 39/45, 6 anomalies via robust outlier rules.
Conditional graph figure with absolute + conditional probabilities.}

\subsection{Introspection: visualisations}
\todoit{shape figure (per-callback typed DAG), process figure
(stream-level DAG), skeleton figures (fig1/2/3), embedding scatter,
skeleton index page. State which scientific question each answers.}

% --- Conclusion -----------------------------------------------------------
\section{Conclusion}
\todoit{summarise: end-to-end stack, validated coverage, novel
shape-skeleton extraction, paper-grade observations on
host-mediated stream serialisation.}

\paragraph*{Threats to validity}
\todoit{(i) heuristic stream-edge inference for regime-A workloads;
(ii) shape-signature granularity is sensitive (54 distinct sigs from 13
callbacks on Autoware); (iii) period extractor not yet implemented,
so frame-level metrics are pending; (iv) tested on x86 + NVIDIA only.}

% --- Future work ----------------------------------------------------------
\section{Future Work}
\todoit{summary of the PhD plan: process-level period extraction
wiring, response-time analysis on heterogeneous taskset, multi-GPU
extension, integration with scheduling decisions.}

% --- Bibliography ---------------------------------------------------------
\bibliographystyle{IEEEtran}
% \bibliography{refs}     % uncomment when you add refs.bib
\begin{thebibliography}{1}
\bibitem{ros2tracing}
\todoit{ros2\_tracing reference}.
\bibitem{nsys}
\todoit{NVIDIA Nsight Systems documentation}.
\bibitem{autoware}
\todoit{Autoware reference}.
\end{thebibliography}

\end{document}
